% 自编教学补充；阅读来源见本章 README。
\begin{frame}[fragile]{1.2A 先选图，再写代码}
\small 类别比较用柱状图；时间变化用折线图；单变量分布用直方图；两个数值的关系用散点图。先用非金融小表练习同一套绘图接口。
\medskip
\begin{lstlisting}
drink_names = ["茶", "咖啡", "果汁"]
drink_counts = [12, 8, 5]
fig, ax = plt.subplots(figsize=(6, 3.5))
ax.bar(drink_names, drink_counts, color="#2B6A6C")
ax.set(title="课堂合成数据：饮品销量", xlabel="饮品", ylabel="销量（杯）")
ax.set_ylim(bottom=0)
fig.tight_layout()
plt.show()
\end{lstlisting}
\end{frame}

\begin{frame}[fragile]{1.2B 折线图：横轴需要有顺序}
\small 连续五日的气温有时间顺序，适合连线；不能为了让曲线好看先按温度大小排序。所有数值均为教学合成数据。
\medskip
\begin{lstlisting}
days = [1, 2, 3, 4, 5]
temperature = [25, 27, 26, 29, 28]
fig, ax = plt.subplots(figsize=(6, 3.5))
ax.plot(days, temperature, marker="o", color="#B40D49")
ax.set(title="连续五日气温（合成）", xlabel="第几日", ylabel="气温（摄氏度）")
ax.set_xticks(days)
fig.tight_layout()
plt.show()
\end{lstlisting}
\end{frame}
\begin{frame}[fragile]{课堂练习：1.2B 折线图：横轴需要有顺序}
在 28 摄氏度画一条虚线并添加图例；为什么折线图不一定要求纵轴从 0 开始？
\medskip
\begin{block}{检查思路}
先写出预期结果，再运行。参考答案见配套 Notebook 的折叠单元格。
\end{block}
\end{frame}
